% ============================================================================
%  SECCIÓN 3.2: IDENTIFICACIÓN DE USUARIOS
% ============================================================================

\section{Identificaci\'on de usuarios}

La identificaci\'on de los usuarios es una de las etapas fundamentales del
an\'alisis de un sistema de software, pues permite delimitar los actores que
interactuar\'an con la soluci\'on y las necesidades que esta debe satisfacer
\citep{sommerville2011, kendall2011}. Para el sistema propuesto se han
identificado tres perfiles de usuario, manteniendo un alcance acotado que
permita desarrollar el producto m\'inimo viable durante la gesti\'on acad\'emica.

\subsection{Perfiles de usuario}

\begin{table}[H]
  \centering
  \caption{Perfiles de usuario identificados}
  \contenidoConFuente{%
    \begin{tablaAPA}
      \begin{tabular}{@{}p{3.0cm}p{3.0cm}p{9.0cm}@{}}
        \toprule
        \textbf{Perfil} & \textbf{Rol} & \textbf{Necesidades principales} \\
        \midrule
        Asistente o aficionado deportivo & Usuario final del sistema &
        Consultar los escenarios cercanos, conocer su ubicaci\'on en el mapa, ver
        las disciplinas, horarios y eventos, y guardar sus escenarios favoritos. \\
        \addlinespace
        Gestor de escenario & Usuario que publica informaci\'on &
        Registrar y actualizar los datos de su escenario (descripci\'on,
        disciplinas, horarios, servicios y eventos) para darle visibilidad. \\
        \addlinespace
        Administrador del sistema & Usuario con privilegios t\'ecnicos &
        Gestionar los usuarios, validar la informaci\'on publicada y administrar
        el cat\'alogo de escenarios y deportes. \\
        \bottomrule
      \end{tabular}
    \end{tablaAPA}%
  }{Elaboraci\'on propia en base a la revisi\'on documental.}
\end{table}
